\section{Delete and Concatenate}
\label{sec:g-delete-concat}
\source{Codeforces 2245B (Notion: Greedy, 1100)}{Medium}

\problemstatement{Array $a$ of $n\le2\cdot10^5$ integers ($|a_i|\le10^9$)
and a cost $c\ge0$. Until the array is empty, either remove one element
(gain its value) or remove two \emph{currently adjacent} elements (gain their
maximum). Every operation costs $c$. Maximise the final score.}

\begin{patternbox}
Start from a baseline (remove everything singly) and ask what each pairing
\emph{changes}: pairing $x\le y$ gains $c-x$ over two single removals.
The best pairs use the smallest elements as their minima.
\end{patternbox}

\subsection*{Approach and intuition}

Removing everything singly gives $\sum a_i-nc$. Replacing two single removals
of $x\le y$ by one pair removal changes the score by
$(\max(x,y)-c)-(x-c)-(y-c)=c-x$. So every pair contributes $c-\min$ of the
pair.

\begin{ideabox}[Which Pairs Are Possible?]
Nested pairs are fine (remove the inner pair first), so any \emph{non-crossing}
matching is achievable. Given $k$ ``small'' elements (the $k$ smallest) and
$k$ ``partners'' (the $k$ largest), a non-crossing matching that pairs every
small element with a partner always exists --- repeatedly match an adjacent
small/partner pair in the order of positions, exactly like balanced
parentheses. Each small element is $\le$ its partner, so the gain is
$\sum_{i\le k}(c-a_{(i)})$. Take pairs while $a_{(i)}<c$ and $i\le n/2$.
\end{ideabox}

\subsection*{Algorithm}
\begin{algorithmic}[1]
\State $score\gets\sum a_i-n\cdot c$; sort $a$
\For{$i=0$ \textbf{while} $i<\lfloor n/2\rfloor$ \textbf{and} $a_i<c$}
  \State $score\gets score+(c-a_i)$
\EndFor
\State \Return $score$
\end{algorithmic}

\dryrun{$a=[3,1,4,1,5,9]$, $c=6$: baseline $23-36=-13$; sorted
$1,1,3,4,5,9$; gains $5,5,3$ for the three smallest (all $<6$). Score $0$.
$a=[1,3,1]$, $c=2$: $5-6=-1$, one pair gains $1$, score $0$. \checkmark}

\subsection*{C++17 implementation}
\cppfile{greedy/20_delete_and_concatenate.cpp}

\begin{complexitybox}
\textbf{Time:} $O(n\log n)$. \textbf{Space:} $O(n)$.
\end{complexitybox}

\begin{pitfallbox}
\begin{itemize}
  \item $n\cdot c$ reaches $2\cdot10^{14}$ and the answer can be
        $-4\cdot10^9$ or lower: \texttt{long long} everywhere.
  \item At most $\lfloor n/2\rfloor$ pairs exist.
\end{itemize}
\end{pitfallbox}
